\documentclass[11pt]{article}
\usepackage{fontspec,amsmath,amssymb,amsthm,hyperref}
\usepackage[margin=1in]{geometry}
\setlength{\emergencystretch}{3em}
\title{A Gaussian-design counterexample using the actual potential Hessian}
\author{LangevinLean mathematical exposition}
\date{}
\begin{document}\maketitle
\noindent Source module: \texttt{AverageLMC/GLMCounterexample.lean}. Compile this proof document with XeLaTeX or LuaLaTeX. The accompanying HTML pairs the mathematical argument with the exact Lean statements and records the scope of the formalization.


This module combines the genuine finite-sum GLM potential and the Gaussian probability obstruction. The design has n=p independent standard Gaussian rows, exactly the isotropic case allowed by the source. The bound event uses the derivative of the actual gradient of the actual potential, so no Hessian identification or concentration estimate is assumed.


\section{Construct precisely the square isotropic Gaussian design}


\texttt{GaussianDesign p} stores p vectors in \(\mathbb R^p\). Its measure is the finite product of p standard Gaussian laws. The product is a probability measure; evaluation at any row preserves the standard Gaussian law, and the coordinate evaluation maps are independent. These facts establish the design assumptions for the counterexample.



\noindent Lean declarations: \texttt{GaussianDesign}, \texttt{gaussianDesignLaw}, \texttt{gaussianDesignLaw\_probability}, \texttt{gaussianDesign\_row\_law}, \texttt{gaussianDesign\_rows\_independent}.

\section{Translate the table into an actual Hessian event}


The isotropic table asserts \(M_2\sqrt p\le Cn\sqrt p\), so at n=p it requires the actual Hessian to be Lipschitz with bound Cp. \texttt{glmHessianBoundEvent} expresses exactly that inequality for all states using \texttt{fderiv} of the actual gradient. The proved rank-one lower bound forces every design in this event to satisfy \(\|X_i\|^3\le Cp\). The row law transfers this event probability to the standard Gaussian measure. For \(p>16\) and \(p>8C^2\), the probability is strictly below one half.



\noindent Lean declarations: \texttt{glmHessianBoundEvent}, \texttt{glmHessianBound\_probability\_lt\_half}.

\section{Rule out every dimension independent table constant}


Assume one nonnegative constant C makes the Hessian bound event have probability at least one half for every positive dimension. Choose a natural p above both 16 and \(8C^2\); the previous theorem gives probability below one half for its first row, a contradiction. The component derivative bound is always the fixed value four, so allowing C to depend only on that bound does not remove the contradiction. The main sampler convergence proofs are independent of this false Gaussian design claim.



\noindent Lean declarations: \texttt{no\_uniform\_glmHessian\_bound}.

\end{document}
